% !TEX program = xelatex
% Собирать XeLaTeX (текст Montserrat, формулы Latin Modern). На Overleaf: Menu -> Compiler -> XeLaTeX
% (или оставьте рядом файл latexmkrc).
\documentclass[aspectratio=169,12pt]{beamer}
\usepackage[lang=ru]{brainpres}

\title[Проксимальный метод]{Негладкая оптимизация. Проксимальный метод}
\subtitle{Методы оптимизации, лекция 5}
\author{Александр Безносиков}
\institute{BRAIn lab, МФТИ}
\date{20 марта 2024}

\begin{document}

\titleframe

\begin{frame}{\contentsname}
  \tableofcontents
\end{frame}

% ------------------------------------------------------------
\section{Постановка задачи}   % слайд раздела вставляется автоматически

\begin{frame}{Задача композитной оптимизации}{Гладкая часть плюс негладкий регуляризатор}
  Рассматриваем задачу
  \[
    \min_{x \in \Rd} \; F(x) = f(x) + r(x),
  \]
  где
  \begin{itemize}
    \item $f$ --- выпуклая и $L$-гладкая;
    \item $r$ --- выпуклая, но, возможно, \alert{негладкая}:
      \begin{itemize}
        \item $r(x) = \lambda \norm{x}_1$ (LASSO);
        \item индикатор выпуклого множества.
      \end{itemize}
  \end{itemize}
\end{frame}

\begin{frame}{Скорости сходимости}
  \begin{columns}[T,onlytextwidth]
    \column{0.31\textwidth}\statcard{$\mathcal{O}(1/k)$}{градиентный спуск на гладкой задаче}
    \column{0.31\textwidth}\statcard{$\mathcal{O}(1/k)$}{проксимальный градиентный метод на композитной задаче}
    \column{0.31\textwidth}\statcard{$\mathcal{O}(1/k^2)$}{с ускорением (FISTA)}
  \end{columns}
\end{frame}

\begin{frame}[c,dark]
  \statement[Главная мысль лекции]{Негладкость не мешает, если у регуляризатора считается prox.}
\end{frame}

% ------------------------------------------------------------
\section{Проксимальный метод}

\begin{frame}{Проксимальный оператор}
  \begin{defn}
    $\operatorname{prox}_{\gamma r}(x) = \argmin_{y} \left\{ r(y) + \frac{1}{2\gamma}\norm{y - x}^2_2 \right\}$.
  \end{defn}
  \begin{thm}
    Если $f$ выпукла и $L$-гладкая, $\gamma \le 1/L$, то
    $F(x^k) - F^* \le \dfrac{\norm{x^0 - x^*}^2}{2\gamma k}$.
  \end{thm}
\end{frame}

\begin{frame}{Три режима}
  \begin{columns}[T,onlytextwidth]
    \column{0.49\textwidth}
      \begin{card}[height=4.6cm]{Градиентный шаг}
        $x^{k+1} = x^k - \gamma \nabla f(x^k)$
        \begin{itemize}
          \item только гладкая часть;
          \item скорость $\mathcal{O}(1/k)$.
        \end{itemize}
      \end{card}
    \column{0.49\textwidth}
      \begin{card}{Проксимальный шаг}
        $x^{k+1} = \operatorname{prox}_{\gamma r}\bigl(x^k - \gamma \nabla f(x^k)\bigr)$
      \end{card}
      \begin{cardgold}{Ускорение}
        FISTA: $\mathcal{O}(1/k^2)$ \cmark
      \end{cardgold}
  \end{columns}
\end{frame}

\begin{frame}{Алгоритм}
  \begin{columns}[c,onlytextwidth]
    \column{0.62\textwidth}
    \begin{algorithm}[H]
      \caption{Проксимальный градиентный метод}
      \begin{algorithmic}[1]
        \Require шаг $\gamma > 0$, точка $x^0$
        \For{$k = 0, 1, \dots, K-1$}
          \State $x^{k+1} = \operatorname{prox}_{\gamma r}\bigl(x^k - \gamma \nabla f(x^k)\bigr)$
        \EndFor
        \Ensure $x^K$
      \end{algorithmic}
    \end{algorithm}
    \column{0.34\textwidth}
    \centering
    \begin{tikzpicture}[node distance=1.2cm]
      \node[presnode] (a) {$x^k$};
      \node[presnode, right=of a] (b) {$y^k$};
      \node[presnode, below=0.8cm of b] (c) {$x^{k+1}$};
      \draw[presedge] (a) -- node[above]{\scriptsize $\nabla f$} (b);
      \draw[presedge] (b) -- node[right]{\scriptsize prox} (c);
    \end{tikzpicture}
  \end{columns}
\end{frame}

\begin{frame}{Важное}
  \begin{columns}[T,onlytextwidth]
    \column{0.48\textwidth}
      \begin{alertblock}{Шаг}
        При $\gamma > 1/L$ метод может расходиться.
      \end{alertblock}
    \column{0.48\textwidth}
      \begin{exampleblock}{Пример}
        Для $r(x) = \lambda\norm{x}_1$ prox --- покомпонентный soft-thresholding.
      \end{exampleblock}
  \end{columns}
\end{frame}

% Шаблоны слайдов с картинкой и цитатой: замените текст на свой
\imageframe{figures/example.jpg}{Слайд с картинкой}{%
  Рисунок или фото занимает левую половину слайда и обрезается под неё
  с сохранением пропорций. Справа --- заголовок и текст.}

\quoteframe{Здесь может быть цитата: она набирается крупно на отдельном слайде.}{Автор цитаты}

% ------------------------------------------------------------
\section{Итоги}

\begin{frame}{Итоги}
  \begin{enumerate}
    \item prox --- обобщение проекции;
    \item скорость как у градиентного спуска: $\mathcal{O}(1/k)$;
    \item ускорение даёт $\mathcal{O}(1/k^2)$.
  \end{enumerate}
  \vfill
  \pdflink{Конспект лекции}{https://example.com/}
\end{frame}

\finalframe{Спасибо за внимание!}{Александр Безносиков \\ BRAIn lab, МФТИ}

\end{document}
